%%%%%%%%%%  scan idx 74  |  folio 61  %%%%%%%%%%
$\Atil$ soumis aux restrictions $\varepsilon(\xi_i)=0$ et $\eta_j=[n_j,1+\eta_j^{(1)}+\cdots+\eta_j^{(nj)}]$} (donc $\eta_j^{(k)}=0$ pour $k>n_j$), \emph{on ait}~:
\[
 R((\lambda^k(\xi_i)),(\lambda^k(\eta_j)))\in\Atil^p
\]
\emph{Soit alors $X$ une variété quasi-projective irréductible non singulière. Quels que soient les éléments $x'_i$ de $\underline K(X)^1$ et les fibrés vectoriels $E'_j$ de dimension $n_j$ sur $X$ ($1\leq i\leq q$, $1\leq j\leq r$), on a}
\[
 R((\lambda^k(x'_i)),(\lambda^k(\gamma(E'_j)))\in\underline K(X)^p\qquad\text{\footnotemark}\ .
\]
\footnotetext{On prouve facilement que, l'hypothèse faite implique si $R$ est exprimé en termes des $\gamma^k(\xi_i)$, $\gamma^k(E_i-\varepsilon(E_i))$, alors tous les termes de $R$ sont de poids $\geq p$ (où $\gamma^k$ est de poids $k$). Alors 2.14 est un cas particulier de X 1.3.2.}

\emph{Démonstration}~: En vertu du \S 2, les $x'_i$ sont eux-mêmes des différences d'éléments de la forme $\gamma(E)$, $E$ étant un fibré vectoriel. D'autre part, en supposant $X$ localement fermé dans un espace projectif $P$, il résulte du \S 2 et de la proposition 2.5 que, pour tout $F\in\cF(X)$, le faisceau $F(n)$ est engendré par ses sections pour $n$ assez grand. Si $F$ est défini par un fibré vectoriel $E$, cela prouve que $E(n)$ est l'image réciproque du fibré vectoriel canonique sur une variété grasmanienne $G$ par un morphisme de $X$ dans $G$. Considérant le morphisme correspondant de $X$ dans $P\times G$, on voit que $E$ est l'image réciproque d'un fibré vectoriel sur $P\times G$. On voit donc, appliquant ceci à tous les fibrés vectoriels qui interviennent, et utilisant le corollaire 2 du théorème 2.12, qu'\emph{on est ramené au cas où} $X$ \emph{est un produit de grasmaniennes}. Dans ce cas, la conclusion résultera du lemme suivant~:

\textbf{Lemme 2.15.}
\emph{Supposons que $X$ soit un produit de grasmaniennes. Alors on peut trouver un $\lambda$-homomorphisme de $\underline K(X)$ dans un $\lambda$-anneau sans torsion de la forme 
